% =====================================================================
\chapter[理解：逆向压缩]{理解：逆向压缩}\label{ch:comprehension}
% =====================================================================

本章处理问题 \ref{prob:mind} 的前半：面对无法从微观通读的大型系统，应当从什么层级
理解，以及如何由理解进入开发。核心命题是理解即逆向压缩，其最优深度由率失真驻点条件
给出，其可行性由观测秩决定。

\section{顺序阅读的失效}\label{sec:why-sequential}

一个自然的阅读方式是追随源代码的物理顺序与控制流：

\begin{quote}
\code{main()} $\to$ \code{module\_a} $\to$ \code{module\_b} $\to$
\code{function\_c} $\to$ \code{helper\_d} $\to$ \code{allocator} $\to$
\code{operating\_system} $\to \cdots$
\end{quote}

然后不断向下追踪，最终陷入无限递归。其失败不是耐心不足，而是结构性的：

\begin{principle}[调用关系不等于认知依赖]\label{pr:call-vs-cog}
函数调用关系是\emph{实现}的拓扑；认知依赖关系是\emph{理解}的拓扑。二者不同构。
研究事务隔离问题时，你可能需要知道事务管理器如何使用锁与快照，但不一定需要理解操作
系统整个虚拟内存子系统；反过来，排查内存使用异常时，事务管理器的高层业务语义可能
暂时不重要。
\end{principle}

更根本的原因是生成方向与理解方向的关系。系统是被\emph{自顶向下}设计的（先有意图，
再有架构，最后才有代码展开）；因此逆向过程应当取同方向。自底向上读代码，等于通过
清点所有分子去理解一只猫：原则上可行，实践上荒谬。

而"原则上可行"这一让步有实际后果：顺序阅读不是\emph{不可能}，而是\emph{被占优}。
命题 \ref{prop:stopping} 给出其精确形式：越过驻点后，描述长度增加
而任务损失不再下降。

\section{认知图}\label{sec:maps}

面对一个大型单体仓库，不应先建立文件地图，而应先建立概念地图。二者的分工是：

\begin{itemize}
\item 文件地图回答"代码在哪里"；
\item 概念地图回答"系统由什么组成、为什么存在这些部分、谁依赖谁、什么规则使它们
      能够协作"。
\end{itemize}

\begin{definition}[四张认知图]\label{def:fourmaps}
对一个系统 $S$ 与任务 $T$，四张认知图分别是：
\begin{enumerate}[label=(\roman*)]
\item \term{系统全景图}：把 $S$ 压缩为 $5$--$15$ 个主要职责域，并回答：最重要的职责域
      是什么？每个职责域为什么存在？它们之间最重要的依赖关系是什么？哪些机制被多个
      职责域共享？哪些状态与行为决定整个系统的正确性？
\item \term{关键子系统图}：对每个关键子系统回答八个问题：职责、边界（负责什么与
      不负责什么）、状态（拥有什么状态、谁有权修改）、机制、接口、约束、失败行为、
      实现入口（哪些源码与测试可以验证上述结论）。
\item \term{关键执行路径图}：把静态架构转换为动态行为模型。例如一条查询路径
      "输入 $\to$ 解析 $\to$ 规划 $\to$ 执行 $\to$ 访问存储 $\to$ 返回"背后的
      真实问题：查询计划何时创建？执行期间状态存放在哪里？事务快照由谁管理？存储
      访问失败时错误如何向上传递？哪些步骤会持有锁？哪些操作可能触发磁盘 I/O？
      查询取消时哪些资源必须释放？
\item \term{契约图}：按定义 \ref{def:relgraph} 的五类边分别记录构件之间的关系，
      特别是每层的\emph{无视集}，即该层被允许永远不关心的内容。
\end{enumerate}
\end{definition}

第 (iv) 张图是本文相对于常规架构文档的补充。它的必要性直接来自命题
\ref{prop:oblivious}：一个没有声明无视集的契约，等于没有声明它的压缩边界，从而
无法判断上层可以依赖它到什么程度。

\section{理解的操作判据}\label{sec:five-capabilities}

"理解"必须被操作化，否则无法判断是否已经理解。本文采用五能力判据。

\begin{definition}[理解判据]\label{def:understand}
称一个智能体（人或机器）\term{理解}了子系统 $M$ 到任务 $T$ 所需的程度，当且仅当它
能够做到：
\begin{enumerate}[label=(U\arabic*)]
\item \term{解释}：说清这个子系统为什么存在；
\item \term{预测}：预测某个输入、状态变化或错误会导致什么行为；
\item \term{定位}：出现故障时判断应该检查哪些组件；
\item \term{修改}：要添加功能时找到正确的扩展点与必须遵守的约束；
\item \term{验证}：知道哪些测试或实现细节能够证明自己的判断。
\end{enumerate}
本文称达到 (U1)--(U5) 的状态为\term{可操作的抽象理解}。
\end{definition}

定义 \ref{def:understand} 与定义 \ref{def:sufficient} 的关系是直接的：(U2) 是
$T$-充分性的行为化表述；(U1)、(U3)、(U4) 是 $T$ 中具体查询族；
(U5) 是"能够追溯到更底层证据"的要求，即定义 \ref{def:mdl} 中"足够准确"与
"能够验证"两条。

一个重要的诊断推论：如果只能复述架构图，却无法预测一个改动会影响什么，那么当前的
抽象还不够。这与第 \ref{sec:seams} 节的缝探测器是同一件事的两个方向。

\section{阅读顺序}\label{sec:protocol}

\begin{principle}[下行协议]\label{pr:protocol}
自顶向下形成模型，按需下钻验证，永远不要顺序阅读。
\end{principle}

具体步骤如下。

\paragraph{学科层。}
读 README、架构文档、架构决策记录、设计论文、内核文档目录。用约 $1\%$ 的时间建立
约 $80\%$ 的骨架模型。产出是一张手绘的器官图：系统有哪几个大块、各自职责、数据
如何流动。

\paragraph{运行观测。}
先跑起来，再读。运行时的行为暴露结构，比静态代码诚实得多：
观察进程与线程模型；打断点或跟踪一次典型请求的完整旅程；看日志与指标，\emph{日志的
词汇表就是系统的自述}。一次端到端的跟踪，抵得上一万行静态阅读。

\paragraph{接口边界。}
按依赖顺序看模块的头文件、公共接口与导出符号，而不是实现。接口是压缩边界
（命题 \ref{prop:oblivious}），它告诉你这一层承诺什么；至于它如何兑现承诺，暂时
不需要。此刻要的是器官图谱，不是组织切片。

\paragraph{数据形状。}
数据模型是最稳定的抽象，控制流最易变。先弄清楚：核心实体是什么、生命周期由谁拥有、
状态机有哪些、谁可写谁只读。握住数据的形状之后，控制流只是数据随时间的重排。

\paragraph{按需下钻。}
只带着具体任务下钻到代码：修这个缺陷、加这个功能。此时阅读是有指向的检索而非漫游。
下钻遵循"调用点 $\to$ 实现 $\to$ 底层机制"的 V 字形：能解释当前问题即收手，然后回到
高层重估模型。

\paragraph{预测与验证。}
在每一层，先写下心智模型（"我猜执行器是这样取下一条元组的"），然后去验证。
\emph{能预测行为就是理解}。预测对了，压缩模型有效；预测错了，那个位置恰好就是系统
最反直觉之处。

\paragraph{回写契约图。}
把新获得的知识写回定义 \ref{def:fourmaps} 的契约图，填入该层的新增保证、新增假设、
以及新增的无视集。这一步是防止认知漂移的唯一手段：若不回写，分层模型会逐渐脱离真实
系统。

\begin{remark}[一个可操作的比例]
深入理解约 $10\%$ 的代码，同时知道另外 $90\%$ 在哪里、归谁管。这个比例不是精确的
目标，而是对"理解不等于穷尽"这一原则的数量化提醒。
\end{remark}

\section{模块边界的探测判据}\label{sec:seams}

架构就是"改动被局部化的那些地方"。因此理解的关键操作是\emph{找缝}。以下判据按可靠性
排序，均可机械执行或半机械执行。

\begin{enumerate}[label=(S\arabic*)]
\item \term{类型停止穿越处}。数据在哪里被重新编码、包装或降级为通用类型
      （如字节数组、字符串、映射）？这是模块边界的强信号：类型系统停止提供保证的
      地方，就是契约必须由人补上的地方。
\item \term{可隔离测试的单位}。哪里存在一个可以在不启动整个系统的前提下跑完的测试
      套件？这是最锋利的判据：\emph{能被隔离验证的单位才是真模块}。它的可靠性来自
      定义 \ref{def:sufficient}：一个单元测试本质上是在断言"接口观察是行为的充分
      统计量"。
\item \term{序列化、版本号与拷贝边界}。哪里有版本字段、序列化、深拷贝、跨进程或
      跨网络调用？跨越这类边界使契约变得可见且可证伪。
\item \term{不变量的集中强制点}。哪里有断言、数据库约束、类型级约束、契约检查？
      不变量被集中强制的地方通常就是层级边界的所在。
\item \term{构建系统的独立目标}。构建配置里哪里存在独立的编译目标或包？这是边界
      的\emph{必要}而非\emph{充分}条件：独立目标未必是真模块，但真模块几乎总有独立
      目标。
\item \term{共同变更图与事故史}。这是最被低估的技术。考察历史提交记录：哪些文件
      倾向于在同一个变更中被一起修改？\emph{共同变更图给出的是语义耦合，而导入图
      给出的是语法耦合}。二者不一致的地方，正是架构文档与现实不符的地方。事故与
      缺陷历史则给出真实的不变量缺口。
\end{enumerate}

\begin{principle}[缝的判据总结]\label{pr:seam-summary}
一个模块是真的，当且仅当它的接口观察是其外部行为的充分统计量；一个边界是真的，
当且仅当跨过它之后，另一侧的变化不再需要被理解。
\end{principle}

第 \ref{ch:design} 章将把原则 \ref{pr:seam-summary} 反过来用作设计目标。

\section{主动辨识与被动阅读}\label{sec:active}

\begin{principle}[理解即系统辨识]\label{pr:sysid}
理解一个复杂系统的本质是\emph{系统辨识}：构造一个能预测其行为的模型，并量化模型的
不确定性。系统辨识的最高效方式是"扰动--观察"循环。
\end{principle}

被动阅读是一种极差的采样方式，因为它\emph{没有误差信号}。阅读产生的每一个判断都
无法被证伪，因此无法被更新。主动探测则相反：先形成一个可证伪的预测（"如果我改这里，
应该只有这两个模块受影响"），做一个最小改动或跑一次测试，然后观察预测误差。
\emph{误差就是你看不见的那部分结构。}

\begin{proposition}[主动探测的优势]\label{prop:active}
设理解过程是模型 $M$ 的迭代更新，每次观测 $o$ 提供关于真实结构的信息
$I(\text{structure}; o \mid M)$。被动阅读的观测分布不依赖于 $M$，因此其期望信息增益
关于 $M$ 是常数；主动探测的观测分布依赖于当前预测 $\hat{y}(M)$，其期望信息增益正比于
预测误差 $\mathbb{E}\,|y - \hat{y}(M)|$。因此在预测误差非零处，主动探测的边际信息
增益严格较高。
\end{proposition}

\begin{proof}[论证]
被动阅读采样的样本与 $M$ 无关，故 $p(o \mid M) = p(o)$，信息增益为常数。
主动探测选择使预测分歧最大的观测（实验设计中的 D-最优或不确定性采样），
$p(o \mid M)$ 集中在 $M$ 当前预测不确定的区域，由信息增益的定义
$I = \mathbb{E}[\log p(o\mid M) - \log p(o)]$，该量非负且当 $p(o\mid M) = p(o)$
时为零，即预测与真实一致时无增益。故在预测误差非零处主动探测严格更优。
\end{proof}

对 AI 而言该结论的分量更重：AI 的优势不在于能"读"更多代码，而在于能高吞吐地\emph{做
探测}，生成假设、跑测试、做最小改动、看共同变更、读事故记录，然后更新模型。这是
人类难以匹敌的频率。

\section{可观测性判据}\label{sec:observability}

第 \ref{sec:comprehension-formal} 节把理解定义为任务约束的压缩，命题
\ref{prop:stopping} 给出下潜的停止点。这两者都预设了"继续观测能改进模型"。本节指出
一个更强的限制：\emph{观测不足时，增加观测量的数量不改变任何东西}。

\begin{definition}[观测算子与任务子空间]\label{def:observability}
设 $\Beh$ 为系统行为空间，$\mathcal{T} \subseteq \Beh$ 为\emph{任务相关差异子空间}，
即定义 \ref{def:task} 中查询族 $Q$ 能够区分的差异所张成的子空间。设
$\mathcal{J} : \Beh \to \mathcal{Y}$ 为智能体\emph{可获得}的观测算子：它把行为映射
到它实际能读取的东西（源码、日志、跟踪、测试结果、运行时指标）。定义
$\mathcal{T}$ 上的\term{观测秩}
\begin{equation}\label{eq:obsrank}
\mathrm{rank}_{\mathcal{T}}\mathcal{J} \;=\;
\dim \mathcal{J}\bigl(\mathcal{T}\bigr),
\end{equation}
即任务相关差异经过 $\mathcal{J}$ 之后仍然可区分的维数。该构造与线性系统理论中的可观测性判据同形\cite{kalman1960}。
\end{definition}

\begin{proposition}[可观测性判据]\label{prop:observability}
智能体能够把系统辨识到 $T$-充分（定义 \ref{def:sufficient}），当且仅当
\begin{equation}\label{eq:obscond}
\mathrm{rank}_{\mathcal{T}}\mathcal{J} \;=\; \dim \mathcal{T}.
\end{equation}
当式 \eqref{eq:obscond} 不成立时，$\mathcal{J}$ 在 $\mathcal{T}$ 上的核非平凡：
存在两次执行 $e_1, e_2$，其全部可获得观测相同而在任务相关行为上不同。此时
\emph{任何数量的额外观测都不能降低任务损失}，因为新增观测仍在同一 $\mathcal{J}$ 的
值域内。
\end{proposition}

\begin{proof}
若式 \eqref{eq:obscond} 成立，则 $\mathcal{J}$ 在 $\mathcal{T}$ 上单射，任务相关差异
可从观测唯一恢复，由定义 \ref{def:sufficient} 即得 $T$-充分性。若不成立，则
$\Ker \mathcal{J} \cap \mathcal{T} \ne \{0\}$，取该核中的非零元对应的一对执行
$e_1 \ne e_2$（在任务相关意义上），二者在 $\mathcal{J}$ 下不可区分，故任何
$\mathcal{J}$-可测的判决规则给出相同答案，至少在其中一个上错误。追加观测不改变
$\mathcal{J}$，故不改变该核。
\end{proof}

命题 \ref{prop:observability} 把"理解不足"分成两种在工程上完全不同的情形：

\begin{principle}[两种理解失败的区分]\label{pr:two-failures}
\begin{enumerate}[label=(\alph*)]
\item \term{秩亏}：观测算子本身不足以区分任务相关差异。此情形的唯一疗法是\emph{改变
      观测}，增加不变量检查、增加跟踪点、增加断言、暴露内部状态。这是设计动作，不是
      阅读动作。
\item \term{秩满但样本不足}：观测足以区分，但尚未采集。此情形的疗法是\emph{增加探测}，
      由命题 \ref{prop:active} 支配其效率。
\end{enumerate}
把二者混淆是"读更多代码"这类建议失效的主要原因：当问题属于 (a) 时，增加阅读量
在信息论上无效。
\end{principle}

\begin{remark}[理解成本的正确计量]\label{rmk:sample-complexity}
情形 (b) 的成本是一个\emph{样本复杂度}，而不是阅读量。若把契约假设限定在一个
假设类 $\mathcal{H}$（定义 \ref{def:operad} 中的运算集合提供了一个自然的候选），
则要以误差 $\varepsilon$ 辨识到 $T$-充分，所需探测次数的下界由 $\mathcal{H}$ 在 $T$
上的区分能力决定，典型形式为 $\Omega(d/\varepsilon)$，其中 $d$ 为该假设类的
容量维数。

这一计量方式给出一个可检验的预测：\emph{设计探测优于穷尽阅读的倍数是可量化的}，
且该倍数随假设类容量的增长而增长。第 \ref{ch:program} 章的假设 $H_5$ 由此获得
一个定量形式。
\end{remark}

\begin{example}[秩亏的典型形态]\label{ex:rank-deficient}
一个服务把全部状态保存在进程内存中，对外只暴露请求的返回值。若任务是判断"并发的两个
请求是否可能读到彼此的部分更新"，则任何数量的返回值观测都不足以判定，因为该差异
完全落在返回值之外。此时正确动作不是读更多代码，而是引入一个可观测量：一个并发不
变量的运行时断言，或一次带交错注入的测试。此例说明原则 \ref{pr:two-failures} 的
区分是实际的而不只是分类学的。
\end{example}

\section{预算与停止条件}\label{sec:budget}

第 \ref{sec:comprehension-formal} 节已经给出停止准则的形式（命题 \ref{prop:stopping}）。
此处给出它的工程化版本与预算分配。

\begin{principle}[下潜预算分配]\label{pr:budget}
把理解预算按"信息增益/成本"排序分配，而不是按代码量或目录顺序分配。典型的高收益
入口依次是：对外接口定义与数据模型 $>$ 端到端执行路径的一次跟踪 $>$ 不变量强制点
$>$ 共同变更热点 $>$ 具体实现细节。低收益入口是：顺序阅读实现体、通读工具函数库、
通读生成代码。
\end{principle}

\begin{principle}[停止条件]\label{pr:stop}
满足下列任一条件即停止下潜：
\begin{enumerate}[label=(\alph*)]
\item 你已能命名管辖当前改动的全部不变量；
\item 继续下潜不再改变你对"这个改动会影响什么"的预测；
\item 下一步的证据获取成本超过了本次改动的预期返工成本。
\end{enumerate}
\end{principle}

条件 (a) 对应原则 \ref{pr:descent}；条件 (b) 对应命题 \ref{prop:stopping}；
条件 (c) 是经济性判断：理解不是越多越好，而是\emph{够用即可}。

\section{机器理解的推论}\label{sec:ai-comprehension}

AI 面临的困难不只是上下文长度有限。即使模型可以读取大量代码，仍然存在三个不同的
问题：信息太多使重要信息被淹没；局部信息之间存在复杂依赖；对细节的理解不一定能自动
形成正确的系统级模型。因此仅仅扩大上下文窗口并不能解决全部问题。

\begin{principle}[AI 的软件认知架构]\label{pr:ai-cog}
一个有效的机器认知架构应当是四层的：
\begin{enumerate}[label=(\roman*)]
\item \term{全局认知模型}：系统目标、架构图、关键不变量、主要数据流、核心设计决策；
\item \term{子系统认知模型}：组件职责、接口、状态所有权、依赖关系、生命周期与错误
      语义；
\item \term{实现级索引}：关键函数、类型定义、调用链、源码位置、测试与运行时证据；
\item \term{验证与反馈}：编译、测试、静态分析、执行跟踪、性能剖析，用于修正上层模型。
\end{enumerate}
其关键性质是：\emph{上层模型不需要复制所有下层信息，只需要保留足够的信息，使智能体
在需要时能够找到正确的下层证据。}
\end{principle}

\begin{definition}[可追溯的子系统认知记录]\label{def:record}
每个子系统的认知记录至少应包含九个字段，见表 \ref{tab:record}。
\end{definition}

\begin{table}[htbp]
\centering
\small
\caption{子系统认知记录的最小字段集}\label{tab:record}
\begin{tabular}{L{3.0cm}L{11.4cm}}
\toprule
\textbf{字段} & \textbf{作用} \\
\midrule
Purpose & 这个子系统为什么存在 \\
Contract & 对外提供什么保证（即 $\Obs_i(\Beh_{i+1})$）\\
Oblivious & 该层被允许永远不关心什么（即 $\Obl(\Lambda_i)$）\\
State & 拥有什么状态，谁有权修改 \\
Dependencies & 依赖哪些其他组件 \\
Invariants & 哪些条件绝不能被破坏，在何处强制 \\
Lifecycle & 关键对象如何创建、运行与销毁 \\
Execution paths & 重要请求或任务如何通过系统 \\
Evidence & 哪些源码、测试或实验支持上述结论 \\
Uncertainty & 哪些内容尚未验证，哪些假设可能错误 \\
\bottomrule
\end{tabular}
\end{table}

表 \ref{tab:record} 中最后两项与 Oblivious 一项是本文的补充，理由如下。

\emph{关于 Uncertainty}：AI 不能把自己生成的一份架构总结当作事实。它需要知道哪些
结论来自实际源码，哪些来自测试，哪些只是暂时的推断。这与第 \ref{ch:objections} 章的
认识论纪律是同一条要求。

\emph{关于 Oblivious}：如上所述，它是压缩边界本身。没有这一项，认知记录只描述了
"有什么"，没有描述"可以不关心什么"，因而无法支撑 (U4) 修改能力。

\emph{关于动态更新}：当代码发生变化时，认知记录必须同步更新。因此有用的机器理解系统
库系统不应只是静态知识库，而应是源码与抽象模型之间\emph{可双向追溯}的关系。

\begin{principle}[可读性从美德变为基础设施]\label{pr:readability}
系统的"理解接口"正在成为一等工程产物。过去文档是给人看的附属品；现在机器可读的
架构描述、接口即文档的实践、以及面向智能体的工作规范，正在成为系统必须设计的组件
，因为代码的读者中多了一个非人类的智能，而它只依赖压缩层。
\end{principle}

本章的经验内容是假设 $H_5$（第 \ref{ch:program} 章）：以分层模型加主动探测为策略的
理解过程，在单位成本上的任务成功率优于穷尽式阅读。命题 \ref{prop:stopping}、
\ref{prop:active} 与 \ref{prop:observability} 在给定模型下成立；定义
\ref{def:understand} 是一个工作定义，其价值取决于它能否被用于区分"理解"与"复述"，
这需要实验检验。
